\section{Introduction}
\label{sec:intro}

Mixture-of-Experts (MoE) architectures~\citep{shazeer2017outrageously,fedus2022switch} have emerged as a scalable approach to conditional computation, routing each input token to a subset of expert sub-networks while keeping total FLOPs manageable. When combined with Vision Transformers (ViTs)~\citep{dosovitskiy2021image,riquelme2021scaling}, MoE layers can increase model capacity without proportionally increasing inference cost. A natural next step is to \emph{prune within} each expert, removing inactive neurons to compress the model further. However, as we show in this work, standard approaches to per-neuron pruning fail when applied inside MoE ViTs.

\paragraph{The gradient-scale problem.}
The dominant strategy for compute-constrained pruning introduces a Lagrangian penalty on FLOPs:
\begin{equation}
\label{eq:lagrangian}
\mathcal{L}_{\text{total}} = \mathcal{L}_{\text{task}} + \lambda\bigl(\text{FLOPs}(\btheta) - B\bigr),
\end{equation}
where $B$ is the target budget and $\lambda$ is a Lagrange multiplier~\citep{gordon2018morphnet,lemaire2019structured}. When each neuron $i$ carries a learnable mask logit $m_i$ passed through a Gumbel-sigmoid gate~\citep{jang2017categorical,maddison2017concrete}, the gradient of the penalty with respect to $m_i$ scales as
\begin{equation}
\label{eq:grad_scale}
\frac{\partial \mathcal{L}_{\text{FLOPs}}}{\partial m_i} \;\approx\; \lambda \cdot \underbrace{\frac{\text{flops}_i}{B}}_{\sim\,10^{-6}} \cdot \underbrace{\frac{\sigma'(m_i/\tau)}{\tau}}_{\sim\,10^{2}} \;\sim\; 10^{-4},
\end{equation}
while task-loss gradients through the same logit are of order $10^{-2}$. The two-order-of-magnitude gap means the Lagrangian signal is too weak to distinguish important neurons from unimportant ones, causing all mask logits to drift uniformly toward $-\infty$ (\emph{uniform collapse}) rather than selectively pruning the least useful neurons (Figure~\ref{fig:motivation}b). This problem is inherent to per-neuron granularity: each individual neuron contributes a negligible fraction of total FLOPs, so the Lagrangian gradient per neuron is vanishingly small.

\begin{figure}[t]
\centering
\includegraphics[width=\linewidth]{figures/fig_motivation.pdf}
\caption{\textbf{Motivation.} (a)~Gradient magnitudes for different parameter groups: mask logits receive $100{\times}$ weaker gradients than task weights. (b)~Under a Lagrangian penalty with shared LR, all mask logits converge to similar values (uniform collapse). (c)~KaleidoNet's cubic schedule with dual-rate optimization produces clear separation between kept and pruned neurons.}
\label{fig:motivation}
\end{figure}

\paragraph{Our approach.}
We propose \textbf{KaleidoNet}, an elastic MoE ViT that sidesteps the Lagrangian gradient-scale problem entirely by replacing the learned penalty with a \emph{deterministic cubic sparsity schedule}~\citep{zhu2017prune}. Instead of relying on gradient-based budget signals, the schedule externally dictates what fraction of neurons should be active at each training step, ramping from 0\% to 70\% sparsity along a cubic curve. Three complementary mechanisms make this schedule effective:

\begin{enumerate}
    \item \textbf{Gradient masking}: Once a neuron is pruned by the schedule, its mask logit is clamped to $-100$ and its gradient is zeroed, preventing accidental recovery and ensuring the pruning decision is irreversible.
    \item \textbf{Dual-rate optimization}: Mask logits are optimized with a $3{\times}$ higher learning rate than task weights (9e-4 vs.\ 3e-4), allowing structural decisions to adapt faster than feature representations.
    \item \textbf{Gumbel-sigmoid straight-through}: During training, mask activations use the Gumbel-sigmoid reparameterization with straight-through gradients~\citep{bengio2013estimating}, bridging the gap between discrete pruning decisions and continuous optimization.
\end{enumerate}

\paragraph{Contributions.} The primary contribution of this paper is diagnostic; KaleidoNet is a practical remedy we use to validate that the diagnosis admits a simple fix.
\begin{enumerate}
    \item \textbf{Primary --- Failure-mode diagnosis.} We identify and formally analyze the \emph{gradient-scale mismatch} that causes Lagrangian FLOPs penalties to fail at per-neuron granularity inside MoE architectures: the per-neuron penalty gradient is ${\sim}100{\times}$ weaker than the task gradient, driving uniform mask collapse rather than selective pruning (Section~\ref{sec:method}).
    \item \textbf{Secondary --- Empirical validation across baselines.} We confirm the diagnosis across three datasets and four pruning baselines: a properly tuned Lagrangian baseline catastrophically collapses to $10.4\%$ on CIFAR-100 (near-random), $36.4\%$ on CIFAR-10, and $4.7\%$ on Tiny-ImageNet---despite reaching the target FLOPs ratio every time. By contrast, \emph{every} schedule-based method we tested (Linear, Random, Magnitude, and our cubic variant) bypasses the gradient-scale issue and reaches near-dense accuracy, isolating the gradient-scale collapse as the cause and the schedule-based formulation as the cure (Section~\ref{sec:experiments}).
    \item \textbf{Tertiary --- KaleidoNet as a stable practical instantiation.} As a concrete remedy we propose KaleidoNet, which replaces the Lagrangian penalty with a deterministic cubic sparsity schedule, gradient masking, and dual-rate optimization. Across all three datasets (50{,}000 steps, 3 seeds), KaleidoNet retains 98.3--99.2\% of dense-model accuracy with $1.8{\times}$ FLOPs reduction and $1.63{\times}$ parameter compression---the strongest schedule-based instantiation we tested, although the schedule shape itself is secondary to the diagnostic insight (Section~\ref{sec:experiments}).
    \item \textbf{Ablation and compression.} Ablations over eight configurations isolate each component's contribution. Model surgery converts soft masks into physically compact networks with $1.63{\times}$ parameter compression (Section~\ref{sec:experiments}).
\end{enumerate}
